\section{Introduction}
In statistical mechanics and computational physics, the \emph{Boltzmann distribution} (also referred to as the \emph{Gibbs distribution} \cite{landau2013statistical}) serves as a fundamental probability measure. It characterizes the likelihood of a system occupying a given state $x$ as a function of the state's energy and the system's temperature. The corresponding probability density function is expressed as
\begin{equation}
    p(x) \propto \exp\left\{-\frac{V(x)}{k_B\mathcal{T}}\right\},
\end{equation}
where $V(x)$ denotes the potential energy of the state $x$, $k_B$ is the Boltzmann constant, and $\mathcal{T}$ is the absolute temperature of the system. The efficient computation and sampling of this Boltzmann distribution remains a central challenge in computational science, underpinning applications across a wide array of disciplines. Specifically, robust sampling facilitates the efficient exploration of high-dimensional phase spaces, which is essential for characterizing rare events \cite{cousins2016reduced, mohamad2016probabilistic, zhang2022koopman}, estimating macroscopic thermodynamic observables \cite{cui2024multilevel, mathias_rousset_free_2010}, and calculating free energy landscapes \cite{darve2001calculating}. Furthermore, these sampling methodologies are deeply integral to advances in Bayesian data assimilation \cite{reich2015probabilistic} and conformational sampling in structural biology \cite{wuttke2014temperature}.

Historically, approaches to this problem have focused on solving the stationary Fokker--Planck equation (SFPE) governing the distribution. Traditional deterministic techniques, such as finite difference or finite element methods \cite{filbet2002numerical, kumar2006solution}, inherently suffer from the \emph{curse of dimensionality}. More recently, the success of deep neural networks has catalyzed the development of machine learning techniques to parameterize high-dimensional PDEs. Pioneering works like the \emph{deep Ritz method} \cite{e_deep_2018}, \emph{physics-informed neural networks} (PINNs) \cite{raissi2019physics}, and methods utilizing weak formulations such as \emph{weak adversarial networks} (WAN) \cite{gu2023stationary, zhai2022deep, lu2024weak, cai2024weak}, have shown promise in approximating the solution of the SFPE. However, these PDE-centric approaches primarily focus on approximating the density function itself. They do not inherently provide a direct, efficient mechanism for generating independent samples from the distribution, which remains a distinct and critical operational requirement in many practical scenarios.

To directly address the sampling requirement and circumvent the dimensionality limitations of grid-based methods, Monte Carlo (MC) methods have established themselves as the dominant paradigm. The standard approach involves numerically integrating the overdamped Langevin stochastic differential equation (SDE) \cite{darve2009computing, ermak_numerical_1980, gilsing2007sdelab}:
\begin{equation}\label{eqn:sdeintro}
    \mathrm{d} X_t = -\nabla V(X_t)\mathrm{d} t + \sqrt{2k_B\mathcal{T}}\mathrm{d} B_t,
\end{equation}
where $B_t$ is a standard Brownian motion. Under the assumption of a strongly confining potential, the probability density of the trajectory $X_t$ converges exponentially fast to the unique invariant Boltzmann distribution \cite{pavliotis_stochastic_2014}. This simplicity and theoretical robustness have made Langevin dynamics the backbone of modern sampling algorithms \cite{cheng_underdamped_2018, dalalyan_sampling_2020, leimkuhler2013rational}.



However, the efficacy of standard Langevin dynamics is severely compromised when the potential landscape $V(x)$ is non-convex and exhibits multiple local minima. In such \emph{metastable} regimes, particularly at low temperatures, the system remains trapped in local basins for exponentially long periods before crossing high energy barriers \cite{freidlin2012random, pavliotis_stochastic_2014}. Consequently, the simulation horizon required to achieve convergence becomes prohibitively large. While advanced techniques such as parallel tempering \cite{lu_methodological_2019} or annealing \cite{neal_annealed_2001} have been proposed to mitigate this, they often introduce significant algorithmic complexity or computational overhead.

Addressing this specific limitation, this work introduces a novel framework that fundamentally accelerates the exploration of complex landscapes. We propose a L\'{e}vy-score-based Monte Carlo method that leverages the heavy-tailed jumps characteristic of L\'{e}vy processes to directly bypass high energy barriers. Crucially, our approach retains the algorithmic simplicity and ease of implementation inherent to standard Monte Carlo methods, while drastically reducing the temporal cost required to escape metastable states.

\smallskip
{\bf Related works:}

\textbf{Fractional Langevin Monte Carlo.} Fractional Langevin Monte Carlo was first introduced in \cite{csimcsekli2017fractional,nguyen2019non}. It has since demonstrated significant utility in modern machine learning, particularly in the domains of non-convex optimization and heavy-tailed sampling \cite{zhu2023distributionally,simsekli2020fractional,ye2018stochastic}. The fundamental principle underlying fractional Langevin Monte Carlo involves formulating a modified SDE of the form $\d X_t=b(X_{t-})\d t + \d L_t^\alpha$. Specifically, the objective is to determine an appropriate drift coefficient $b$ such that the resulting dynamics, driven by an $\alpha$-stable L\'{e}vy process $L_t^\alpha$, preserve the target invariant measure associated with \eqref{eqn:sdeintro}. However, these foundational studies omit rigorous mathematical arguments demonstrating that the considered target measures are genuinely invariant under the proposed dynamics. Furthermore, they lack a formal proof regarding the uniqueness of the stationary solution to Kolmogorov's forward equation, relying instead on heuristic justifications rather than complete theoretical guarantees. Subsequent research has sought to generalize this framework. For instance, the authors of \cite{oechsler2024levy} extended the methodology to encompass general L\'{e}vy processes; however, their analytical scope was restricted to the specific setting of target distributions supported on the half-line and driven by compound Poisson noise. Additionally, recent studies \cite{zhang2023ergodicity,behme2025evy} have adapted the approach to investigate specific target invariant distributions exhibiting heavy-tailed properties.



\textbf{Score-based generative models with $\alpha$-stable L\'{e}vy processes.}
In \cite{yoon2023score}, a class of score-based generative models was introduced, employing a tractable forward L\'{e}vy process defined by $X_t = a(t)X_0 + \gamma(t)\epsilon$. In this formulation, the initial state $X_0$ is drawn from the target data distribution, and the driving noise $\epsilon$ follows an isotropic $\alpha$-stable distribution with a unit scale parameter. To rigorously construct the corresponding time-reversed L\'{e}vy dynamics, the authors define a fractional score function, $s_t^{(\alpha)}(x) = (\triangle^{\frac{\alpha-2}{2}} \nabla p_t(x))/p_t(x)$ for $1 < \alpha \le 2$, associated with the marginal probability density $p_t(x)$. Computationally, this fractional score function is parameterized via a neural network. Its analytical structure is amenable to efficient approximation using denoising score matching, thereby closely paralleling the established training paradigms of standard score-based diffusion models.


\subsection*{Our Contributions}
The primary objective of this work is to overcome the profound computational bottlenecks associated with sampling from metastable invariant distributions. Building upon our previous theoretical investigations, we develop a novel, mathematically rigorous L\'{e}vy-score-based Monte Carlo sampling framework. The main contributions of this paper are summarized as follows:

\begin{itemize}
    \item \textbf{Robust and efficient method:} We propose a fundamentally new sampling dynamics by systematically injecting jump noise into the standard continuous It\^{o} diffusion. Specifically, we construct a modified drift function that explicitly incorporates a L\'{e}vy score. This carefully engineered drift balances the introduced heavy-tailed jumps, ensuring that the augmented stochastic system still exactly preserves the target Boltzmann distribution as its invariant measure.
    \begin{enumerate}
        \item \textbf{Accelerated transitions across metastable states:} We effectively resolve the critical challenge of slow mixing in non-convex potential landscapes. The explicitly introduced jump mechanisms allow the system's trajectory to rapidly bypass high-energy barriers. This significantly accelerates transitions between isolated metastable states, fundamentally mitigating the exponentially long trapping times that plague classical continuous-path Langevin dynamics.
    
    \item \textbf{Algorithmic simplicity and empirical efficacy:} Despite the sophisticated theoretical machinery underlying the SDE construction, the resulting Monte Carlo scheme retains the exceptional algorithmic simplicity and ease of implementation characteristic of classical methods. Through extensive numerical experiments on benchmark problems exhibiting complex metastability, we empirically demonstrate the superior temporal efficiency, scalability, and robustness of our approach.
    \end{enumerate}
    
    \item \textbf{Rigorous theoretical guarantees:} A cornerstone of the proposed methodology lies in its rigorous mathematical foundation. For the newly constructed jump-augmented SDE, under some mild assumptions, we establish its global well-posedness and geometric ergodicity, alongside a detailed theoretical analysis of the mechanisms driving the accelerated sampling rate. These theoretical guarantees are crucial, as they ensure that the time-averaged trajectory converges robustly to the target distribution, thereby providing a strict justification for the proposed Monte Carlo sampling algorithm (see Theorems \ref{theo:well-posedness}, \ref{theo:p_infty}, and \ref{thm:gap_from_shell_cheeger}).
    
    \item \textbf{Numerical verification:} We conduct numerical experiments on both low and highdimensional problems, with the presence of meta-stable states. We compare
our LSB-MC method with the traditional It\^{o} diffusion-based Monte Carlo and demonstrate that the LSB-MC achieves comparable
accuracy with a significantly lower computational cost. Furthermore, the LSB-MC can explore
all the meta-stable states in both low temperature and high temperature scenarios
\end{itemize}

The remainder of this paper is organized as follows. Section \ref{sec:preliminaries} provides the preliminary mathematical background on It\^{o} diffusions and the target Boltzmann distribution. In Section \ref{sec:method}, we detail the construction of the jump-noise-augmented SDE, and the rigorous proofs of global well-posedness and ergodicity, and also the mechanism of acceleration is discussed. Section \ref{sec:experiments} presents comprehensive numerical experiments to validate the efficacy of the proposed Monte Carlo method. Finally, concluding remarks are given in Section \ref{sec:conclusion}.


